\chapter{Background}
\label{chap:background}

This chapter provides the conceptual foundations needed for the rest of the thesis. We begin with a brief introduction to large language models (Section~\ref{sec:bg-llms}) and to personally identifiable information and its detection (Section~\ref{sec:bg-pii}), kept deliberately concise as they serve only as context. The bulk of the chapter (Section~\ref{sec:bg-zkp}) develops zero-knowledge proofs in depth, since they are the cryptographic core of this work. We close with a description of EzKL and zero-knowledge machine learning (Section~\ref{sec:bg-ezkl}), the practical toolchain on which our implementation is built.

\section{Large Language Models}
\label{sec:bg-llms}

\begin{sloppypar}
A large language model (LLM) is a neural network trained on large corpora of text to model the probability distribution of natural language. Modern LLMs are based on the \emph{Transformer} architecture~\cite{vaswani2017attention}, which uses self-attention mechanisms to capture dependencies between tokens regardless of their distance in the input sequence. At a high level, an LLM is an autoregressive predictor: given a sequence of tokens, it produces a probability distribution over the next token, and text is generated by repeatedly sampling from this distribution and appending the result to the context.
\end{sloppypar}

\medskip

\begin{sloppypar}
For the purposes of this thesis, the internal mechanics of the model are not essential. What matters is a single observation: the output of an LLM is unconstrained natural-language text, produced from patterns learned during training. Because that training data may itself contain personal information, and because users may supply personal information in their prompts, an LLM can reproduce or generate text containing personally identifiable information in its responses. This is true regardless of the model's size or quality. Consequently, any deployment that handles personal data must treat the model output as untrusted from a privacy standpoint and apply a filtering step before the output is delivered or acted upon. In our reference implementation we use DistilGPT-2~\cite{distilgpt2}, a compact version of GPT-2 obtained through knowledge distillation~\cite{sanh2019distilbert}, as a representative generator; the attestation mechanism we develop is independent of the specific model used.
\end{sloppypar}

\section{Personally Identifiable Information and its Detection}
\label{sec:bg-pii}

Personally identifiable information (PII) is any data that can be used to identify a specific individual, either on its own or in combination with other information. Common examples include names, email addresses, telephone numbers, postal addresses, national identification numbers, and financial account numbers. Under the European General Data Protection Regulation, such data falls within the broader category of \emph{personal data}, and its processing is subject to strict legal obligations, which we examine in detail in Chapter~\ref{chap:related-work}.

\medskip

In practice, PII detection in text is performed using one of two broad families of techniques, often in combination:

\medskip

\textbf{Rule-based detection.} Many forms of PII follow well-defined structural patterns. Email addresses, phone numbers, credit-card numbers, and identification numbers can be matched with regular expressions (regex), which describe a search pattern as a sequence of literal and meta characters. Rule-based detection is fast, deterministic, and interpretable: a given input either matches a pattern or it does not. Its weakness is that it only captures information with predictable structure, and it cannot recognise, for example, that an arbitrary string of words constitutes a person's name.

\medskip

\textbf{Model-based detection.} To capture unstructured PII such as personal names, organisations, or locations, systems employ \emph{named entity recognition} (NER), a sequence-labelling task in which a model assigns each token a label indicating whether it forms part of an entity of interest. Modern NER systems are typically built on top of pretrained Transformer encoders such as BERT~\cite{devlin2019bert}, fine-tuned on annotated entity data. Model-based detection is more flexible than regex but is also probabilistic, heavier to run, and harder to reason about formally.

\medskip

The system described in this thesis uses a regex-based detector as its operational filter and, as explained in Chapter~\ref{chap:system-design}, a deliberately simplified byte-level policy as the computation that is actually attested in zero knowledge. The reasons for this separation, and an initial attempt to attest a full NER model, are discussed there.

\section{Zero-Knowledge Proofs}
\label{sec:bg-zkp}

A zero-knowledge proof (ZKP) is a protocol that allows one party, the \emph{prover}, to convince another party, the \emph{verifier}, that a given statement is true, while revealing nothing beyond the validity of the statement itself. The concept was introduced by Goldwasser, Micali, and Rackoff~\cite{goldwasser1989knowledge}, who formalised the notion of the ``knowledge complexity'' of an interactive proof, that is, the amount of knowledge the verifier gains beyond the truth of the assertion. In a zero-knowledge proof, that amount is, informally, zero.

\medskip

To make this precise, it is useful to recall the underlying computational model. A decision problem can be framed as a \emph{language} $L$, the set of all strings (inputs) for which the answer is ``yes''. Proving the statement ``$x \in L$'' then means convincing the verifier that the input $x$ belongs to the language. For the statements relevant to this thesis, membership corresponds to the existence of a \emph{witness} $w$, a private value, such that a known relation $R(x, w)$ holds; for example, ``there exists an input that, when fed to a certified filter, produces a \texttt{PASS} result''. The challenge addressed by zero-knowledge proofs is to convince the verifier that such a $w$ exists without revealing $w$ itself.

\subsection{Interactive Proof Systems}
\label{subsec:bg-interactive}

The original formulation of zero-knowledge proofs is \emph{interactive}: prover and verifier exchange a sequence of messages, and the verifier ultimately decides whether to accept or reject. In a typical round, the verifier issues a random \emph{challenge} and the prover must respond with an answer consistent with the claimed statement. The use of randomness is essential: because the prover cannot predict the challenge in advance, it cannot prepare a fraudulent response ahead of time, and only a prover that genuinely possesses the witness can answer correctly with high probability across many rounds. With each additional round, the probability that a cheating prover escapes detection decreases exponentially.

\medskip

This interaction structure can be modelled as two communicating Turing machines that alternate activity, one computing while the other waits, until the verifier reaches a decision. Formally, the framework is captured as follows.

\begin{definition}[Interactive proof system]
\label{def:interactive-proof}
Let $L$ be a language. An \emph{interactive proof system} for $L$ is a pair $(P, V)$, where $V$ is a probabilistic polynomial-time interactive machine (the verifier) and $P$ is an interactive machine of unbounded computational power (the prover). On common input $x$, the two machines exchange a sequence of messages, after which $V$ outputs a verdict in $\{\mathtt{accept}, \mathtt{reject}\}$. We write $\langle P, V \rangle(x)$ for the random variable denoting $V$'s output at the end of the interaction on input $x$.
\end{definition}

The interactive setting is powerful and forms the theoretical basis for the field, but it has a practical drawback: the prover must produce a separate, live interaction for each verifier, and both parties must be online simultaneously. We return to how this limitation is overcome in Section~\ref{subsec:bg-noninteractive}.

\subsection{The Three Defining Properties}
\label{subsec:bg-properties}

A zero-knowledge proof system is characterised by three properties.

\medskip

\textbf{Completeness.} If the statement is true and both parties follow the protocol honestly, the verifier accepts with high probability. In other words, an honest prover holding a valid witness can always convince an honest verifier.

\medskip

\textbf{Soundness.} If the statement is false, no cheating prover can convince the verifier to accept, except with negligible probability. This is what protects the verifier: a prover cannot prove something untrue.

\medskip

\textbf{Zero-knowledge.} The verifier learns nothing from the interaction beyond the fact that the statement is true. In particular, it gains no information about the witness $w$.

\medskip

Completeness and soundness together are what make the protocol a \emph{proof}; the zero-knowledge property is what makes it \emph{private}. The first two are shared with ordinary interactive proofs, while the third is the distinctive contribution of the zero-knowledge model. The first two can be stated precisely as follows.

\begin{definition}[Completeness and soundness]
\label{def:completeness-soundness}
Let $(P, V)$ be an interactive proof system for a language $L$, and let $c, s : \N \to [0,1]$ be functions with $c(|x|) > s(|x|)$. The system satisfies:
\begin{itemize}
    \item \textbf{Completeness} with parameter $c$ if, for every $x \in L$,
    \[
        \Pr\!\left[\, \langle P, V \rangle(x) = \mathtt{accept} \,\right] \;\ge\; c(|x|);
    \]
    \item \textbf{Soundness} with parameter $s$ if, for every $x \notin L$ and every (possibly cheating) prover strategy $P^{*}$,
    \[
        \Pr\!\left[\, \langle P^{*}, V \rangle(x) = \mathtt{accept} \,\right] \;\le\; s(|x|).
    \]
\end{itemize}
The quantity $s(\cdot)$ is the \emph{soundness error}. In the ideal case $c = 1$ and $s$ is a negligible function; in practice the soundness error is driven below any desired threshold by repeating the protocol, since independent repetitions cause it to decay exponentially.
\end{definition}

\subsection{Formalising Zero-Knowledge: the Simulation Paradigm}
\label{subsec:bg-simulation}

The zero-knowledge property is defined using a construction known as the \emph{simulator}. Intuitively, we want to say that the verifier ``learns nothing'' from the proof. This is formalised by requiring that whatever the verifier sees during a real interaction with the prover could have been produced by the verifier alone, without any access to the prover or the witness.

\medskip

Concretely, a \emph{simulator} is an algorithm that, given only the public input $x$, produces a \emph{transcript}, a complete record of the messages that would be exchanged, that is indistinguishable from the transcript of a genuine interaction between the prover and the verifier. If such a simulator exists, then the real interaction cannot have conveyed any knowledge: anything the verifier could compute after talking to the prover, it could equally have computed on its own by running the simulator. A protocol is therefore zero-knowledge if, for every possible (including dishonest) verifier $V^{*}$, there exists a simulator $S$ such that the distribution of simulated transcripts is indistinguishable from the distribution of real transcripts.

\medskip

The strength of the guarantee depends on the notion of indistinguishability used:

\begin{itemize}
    \item \textbf{Perfect zero-knowledge:} the simulated and real transcripts are identically distributed.
    \item \textbf{Statistical zero-knowledge:} the two distributions are statistically close, differing by a negligible amount.
    \item \textbf{Computational zero-knowledge:} the two distributions cannot be told apart by any efficient (probabilistic polynomial-time) algorithm.
\end{itemize}

Most practical proof systems, including the ones used in this work, achieve computational zero-knowledge, which is stated formally below.

\begin{definition}[Computational zero-knowledge]
\label{def:czk}
An interactive proof system $(P, V)$ for a language $L$ is \emph{computational zero-knowledge} if, for every probabilistic polynomial-time verifier $V^{*}$, there exists a probabilistic polynomial-time algorithm $S$ (the \emph{simulator}) such that the two ensembles
\[
    \big\{ \langle P, V^{*} \rangle(x) \big\}_{x \in L}
    \qquad \text{and} \qquad
    \big\{ S(x) \big\}_{x \in L}
\]
are computationally indistinguishable; that is, no probabilistic polynomial-time algorithm can distinguish them with non-negligible advantage. The output of $\langle P, V^{*} \rangle(x)$ here denotes the full transcript of the interaction together with $V^{*}$'s internal state.
\end{definition}

Intuitively, the definition says that whatever a verifier learns from interacting with the prover, it could have produced by itself by running the simulator $S$ on the public input alone; hence the interaction conveys no additional knowledge. The privacy of such systems thus rests on standard cryptographic hardness assumptions, namely that no feasible adversary can distinguish the two ensembles. A subtle but important point is that the guarantee must hold even when the verifier already possesses some side information related to the secret before the protocol begins; a robust definition of zero-knowledge accounts for such \emph{auxiliary inputs}, ensuring the proof leaks nothing even to a verifier that is not starting from scratch.

\subsection{Non-Interactive Proofs and the Fiat--Shamir Transform}
\label{subsec:bg-noninteractive}

The requirement that prover and verifier interact in real time is a serious obstacle for many applications, including ours, where we would like to generate a proof once and have it checked later by any number of parties. A \emph{non-interactive} proof removes this requirement: the prover produces a single proof object that any verifier can check on its own, with no further communication. This naturally supports \emph{offline verification}, in which the verifier is not present at the time the proof is generated and simply receives the proof later, for instance through a public ledger or alongside a delivered response.

\medskip

The standard tool for converting an interactive protocol into a non-interactive one is the \emph{Fiat--Shamir transform}~\cite{fiat1986how}. The key insight is that the verifier's role in an interactive protocol is often limited to supplying unpredictable random challenges. The Fiat--Shamir transform replaces these random challenges with the output of a cryptographic hash function applied to the protocol transcript so far. Because a secure hash function behaves, for practical purposes, like an unpredictable random oracle, the prover cannot anticipate or manipulate the resulting ``challenge'' any more than it could a genuinely random one. The prover can thus compute the entire transcript by itself, producing a self-contained proof that requires no live verifier.

\subsection{From Computation to Circuits}
\label{subsec:bg-circuits}

Zero-knowledge proofs operate on mathematical statements, not on arbitrary programs. To prove a statement about a computation, that computation must first be expressed in a form the proof system can handle, typically an \emph{arithmetic circuit} or, equivalently, a system of constraints such as a \emph{Rank-1 Constraint System} (R1CS). In this representation, the entire computation, including every intermediate value, is encoded as a set of equations over a finite field that the prover's witness must satisfy.

\medskip

This translation step, often called \emph{arithmetisation}, is central to applying zero-knowledge proofs in practice. The statement ``the program $P$ on private input $x$ produces output $o$'' becomes ``there exists an assignment to the circuit's variables, consistent with the private input $x$, that satisfies every constraint and yields $o$''. The cost of proving, and the size of the resulting circuit, depend directly on the complexity of the computation being arithmetised. This observation becomes important in Chapter~\ref{chap:system-design}, where the cost of expressing a neural network as a circuit drives a central design decision.

\subsection{Commitment Schemes}
\label{subsec:bg-commitments}

A \emph{commitment scheme} is a cryptographic primitive that behaves like a sealed envelope. It allows a party to commit to a chosen value while keeping it hidden, and to reveal it later in a way that cannot be altered after the fact.

\begin{definition}[Commitment scheme]
\label{def:commitment}
A \emph{commitment scheme} is a pair of algorithms $(\mathsf{Commit}, \mathsf{Open})$. The committer computes $c = \mathsf{Commit}(m, r)$ on a message $m$ and randomness $r$, publishes $c$, and later reveals $(m, r)$ so that any party can check $c = \mathsf{Commit}(m, r)$ via $\mathsf{Open}$. The scheme must satisfy:
\begin{itemize}
    \item \textbf{Hiding:} the commitment $c$ reveals no feasibly computable information about $m$; formally, for any two messages $m_0, m_1$, the distributions of $\mathsf{Commit}(m_0, r)$ and $\mathsf{Commit}(m_1, r)$ over random $r$ are (computationally or statistically) indistinguishable.
    \item \textbf{Binding:} it is infeasible to find $(m, r) \neq (m', r')$ with $m \neq m'$ such that $\mathsf{Commit}(m, r) = \mathsf{Commit}(m', r')$; once published, $c$ cannot be opened to two different messages.
\end{itemize}
\end{definition}

Informally, one can place a value inside the envelope and seal it; the value can later be revealed, but it cannot be swapped for a different one in the meantime.

\medskip

Commitments are pervasive in zero-knowledge constructions, where they are used to bind the prover to its inputs and intermediate values without revealing them. In this thesis, a cryptographic hash function serves as a commitment to the attested model output: the hash is published openly and binds the proof to a specific output, while revealing nothing about the output's contents. The precise role and limitations of this commitment are discussed in Chapter~\ref{chap:system-design}.

\subsection{zk-SNARKs and zk-STARKs}
\label{subsec:bg-snark-stark}

Two families of non-interactive zero-knowledge proof systems dominate current practice. They share the same goals but differ in their cryptographic foundations and resulting trade-offs.

\medskip

\textbf{zk-SNARKs} (Zero-Knowledge Succinct Non-interactive Arguments of Knowledge) are characterised above all by their \emph{succinctness}: proofs are small and verification is very fast, typically milliseconds, regardless of the size of the underlying computation. These properties have made zk-SNARKs the dominant choice for practical zero-knowledge applications. Their cryptography is generally based on elliptic-curve pairings. The classical drawback is that many zk-SNARK constructions require a \emph{trusted setup} (discussed below), and in the most basic schemes this setup must be repeated for every distinct circuit. Construction families such as Groth16~\cite{groth2016size} achieve very compact proofs of a few hundred bytes, while PLONK~\cite{gabizon2019plonk} introduced a \emph{universal} setup that can be reused across circuits, at the cost of somewhat larger proofs; universal-setup systems based on Halo2 and KZG commitments, such as the EzKL backend used in this thesis, produce proofs on the order of tens of kilobytes. A practical caveat of pairing-based zk-SNARKs is that elliptic-curve cryptography is not, with current knowledge, secure against quantum adversaries.

\medskip

\textbf{zk-STARKs} (Zero-Knowledge Scalable Transparent Arguments of Knowledge) take a different approach. They require \emph{no trusted setup}, deriving all necessary parameters from public randomness, a property referred to as \emph{transparency}. Their security rests on the collision-resistance of hash functions rather than on elliptic curves, which makes them plausibly \emph{post-quantum} secure. Internally, they express the computation trace using Merkle trees and rely on low-degree testing protocols such as FRI to let the verifier challenge the prover at a small number of random points. The price for these advantages is proof size: zk-STARK proofs are substantially larger than zk-SNARK proofs, ranging from tens to hundreds of kilobytes, and they are more computationally demanding to generate.

\medskip

Table~\ref{tab:snark-stark} summarises the principal trade-offs. The system in this thesis uses a zk-SNARK, both because succinct proofs and fast verification are well suited to an attestation that may be checked many times by auditors, and because the EzKL toolchain we build on targets this family.

\begin{table}[h]
\centering
\caption{Comparison of zk-SNARKs and zk-STARKs (representative figures; zk-SNARK proof size is construction-dependent, see text).}
\label{tab:snark-stark}
\begin{tabular}{l l l}
\hline
\hline
Property & zk-SNARK & zk-STARK \\
\hline
Proof size & Small (hundreds of B to tens of KB) & Large (tens--hundreds of KB) \\
Verification time & Milliseconds & Milliseconds (slightly higher) \\
Trusted setup & Often required & Not required (transparent) \\
Cryptographic basis & Elliptic-curve pairings & Hash functions \\
Post-quantum security & No & Plausibly yes \\
\hline
\hline
\end{tabular}
\end{table}

\subsection{Trusted Setup}
\label{subsec:bg-setup}

Many zk-SNARK systems require a one-time \emph{trusted setup} phase that produces a set of public parameters, sometimes called a structured reference string (SRS), used by both prover and verifier. The setup involves secret randomness that must be discarded afterwards: if an adversary were to learn this secret, often referred to as ``toxic waste'', it could forge proofs of false statements, breaking soundness. To mitigate this risk, production systems generate the SRS through a \emph{ceremony} distributed across many independent participants, such that the setup remains secure as long as at least one participant behaves honestly and destroys their portion of the secret.

\medskip

The distinction between \emph{circuit-specific} and \emph{universal} setups is practically important. A circuit-specific setup must be regenerated whenever the circuit changes, whereas a universal setup (as in PLONK) can be reused for any circuit up to a bounded size. This distinction has direct consequences for a system like ours, in which the attested policy is encoded as a fixed circuit: any change to the policy requires regenerating the setup, a limitation we revisit in Chapter~\ref{chap:discussion}. For the experiments in this thesis, the SRS is generated locally, which is appropriate for a research prototype but, as discussed later, would be replaced by a public ceremony in a real deployment.

\section{Zero-Knowledge Machine Learning and EzKL}
\label{sec:bg-ezkl}

\emph{Zero-knowledge machine learning} (ZKML) is the application of zero-knowledge proofs to machine-learning computations: proving that a model produced a particular output on a given input, or that an inference was carried out faithfully, without necessarily revealing the input, the model parameters, or both. The central difficulty is the arithmetisation cost discussed in Section~\ref{subsec:bg-circuits}: neural networks involve large numbers of multiplications and non-linear operations, all of which must be expressed as circuit constraints, and the resulting circuits can be very large. A growing body of work has tackled this challenge, both through specialised constructions for particular model classes~\cite{liu2021zkcnn,lee2024vcnn} and through general-purpose compilers that turn model descriptions into provable circuits.

\medskip

\textbf{EzKL}~\cite{ezkl} is a toolchain of the latter kind, and the one used throughout this thesis. EzKL takes a machine-learning model exported in the open ONNX (Open Neural Network Exchange) format, compiles it into an arithmetic circuit, and provides commands to generate the setup parameters, produce zk-SNARK proofs over the model's inference, and verify them. This pipeline makes it possible to attest the execution of a model without hand-crafting a circuit, which is precisely what our system requires: we express a filtering policy as a small model, export it to ONNX, and let EzKL compile and prove it.

\medskip

EzKL's reliance on ONNX and on its underlying inference backend imposes practical constraints, however. Not every operation that can appear in an ONNX graph is supported by the proving backend, and large models can exhaust available memory during the compilation or setup phases. These constraints are not incidental; as Chapter~\ref{chap:system-design} describes, they directly shaped the design of the attested policy in our implementation, forcing a move away from a full neural detector towards a lightweight, circuit-friendly alternative.
